\chapter{مبانی نظری و فناوری‌های مورد استفاده}

\section{پایش عملکرد و مشاهده‌پذیری}
مشاهده‌پذیری در نرم‌افزار به توانایی فهم وضعیت داخلی سامانه بر اساس خروجی‌های قابل مشاهده مانند رخدادها، معیارها و گزارش‌ها گفته می‌شود. در سامانه‌های مبتنی بر \lr{JVM}، رخدادهای عملکردی می‌توانند از درون برنامه، ماشین مجازی، کتابخانه‌های زمان اجرا و سیستم‌عامل به دست آیند. ترکیب این منابع دید کامل‌تری نسبت به رفتار برنامه ایجاد می‌کند.

\section{\lr{Java Flight Recorder}}
\lr{JFR} بخشی از سکوی جاوا است که رخدادهایی درباره اجرای برنامه، مصرف پردازنده، حافظه، جمع‌آوری زباله، نخ‌ها، قفل‌ها، ورودی/خروجی و سایر رفتارهای زمان اجرا تولید می‌کند \cite{jfrdocs}. مزیت اصلی \lr{JFR} این است که برای پایش کم‌سربار طراحی شده و داده‌های آن با ساختار مشخص قابل مصرف است. در پروژه مورد بررسی، ماژول \lr{event-receiver} به فرایند هدف متصل شده و رخدادهای \lr{JFR} را به پیام‌های قابل ارسال در \lr{Kafka} تبدیل می‌کند.

\section{\lr{eBPF} و رخدادهای سطح هسته}
\lr{eBPF} امکان اجرای برنامه‌های کوچک و کنترل‌شده در فضای هسته لینوکس را فراهم می‌کند و برای پایش شبکه، فراخوانی‌های سیستمی، امنیت و عملکرد کاربرد زیادی دارد \cite{ebpfdocs}. در این پروژه، بخش \lr{kernel-event-sender} با زبان \lr{C} و \lr{libbpf} رخدادهای سیستمی را جمع‌آوری می‌کند و بخش \lr{kernel-event-receiver} آن‌ها را دریافت و به جریان پیام تبدیل می‌کند. این مسیر باعث می‌شود سامانه علاوه بر رخدادهای سطح \lr{JVM}، نشانه‌های نزدیک‌تر به سیستم‌عامل را نیز ببیند.

\section{\lr{Kafka} و پردازش جریانی}
\lr{Kafka} یک سکوی پیام‌رسانی توزیع‌شده برای انتقال جریان داده است و \lr{Kafka Streams} کتابخانه‌ای برای پردازش جریانی روی داده‌های ذخیره‌شده در \lr{Kafka} محسوب می‌شود \cite{kafkadocs,kafkastreams}. در این پروژه، رخدادهای خام در موضوع‌هایی مانند \lr{\texttt{profiler.raw.events}} و \lr{\texttt{profiler.kernel.raw.events}} قرار می‌گیرند. سپس ماژول \lr{event-aggregator} این رخدادها را در پنجره‌های زمانی ده‌ثانیه‌ای تجمیع و خروجی را در موضوع‌های تجمیع‌شده منتشر می‌کند.

\section{تشخیص ناهنجاری}
تشخیص ناهنجاری به شناسایی رفتارهایی گفته می‌شود که با الگوی معمول سامانه تفاوت معنادار دارند. در این پروژه، رویکردهای ساده اما قابل توضیح استفاده شده‌اند: قواعد آستانه‌ای برای معیارهایی که حد قابل قبول مشخص دارند، خط پایه غلتان برای مقایسه مقدار فعلی با گذشته نزدیک، و تشخیص روند افزایشی برای معیارهایی که رشد پیوسته آن‌ها می‌تواند نشانه مشکل باشد. این انتخاب‌ها برای یک پروژه مهندسی مناسب هستند، زیرا خروجی آن‌ها قابل تفسیر و پیاده‌سازی آن‌ها کم‌هزینه است.

\section{جمع‌بندی}
ترکیب \lr{JFR}، \lr{eBPF} و \lr{Kafka Streams} امکان ایجاد خط لوله‌ای را فراهم می‌کند که هم رخدادهای درون ماشین مجازی و هم رخدادهای سیستم‌عامل را دریافت کند. این فصل مفاهیم لازم برای فهم طراحی سامانه را معرفی کرد و فصل بعد معماری پیشنهادی را توضیح می‌دهد.
